\documentclass[10pt,a4paper]{amsart}
\usepackage[margin=19mm]{geometry}
\usepackage{amsmath,amssymb,mathtools}
\usepackage[scr=rsfs,cal=euler]{mathalfa}
\usepackage{xfrac}
\usepackage[unicode,hidelinks,bookmarks=false]{hyperref}
\newcommand{\ie}{i.e.\ }
\newcommand{\cf}{cf.\ }
\newcommand{\resp}{respectively\ }
\newcommand{\Subsection}{\subsection}
\providecommand{\nobreakdash}{\nobreak\hbox{-}}
\setlength{\emergencystretch}{2em}
\multlinegap0pt
\usepackage{enumerate}
\usepackage{amssymb}
\newtheorem{corollary}{Corollary}
\newtheorem{lemma}{Lemma}
\title{\bf A Sharp Sobolev Inequality on Compact CR Manifolds}
\author{Zhongwei Tang, Bingwei Zhang}
\date{Source: arXiv:2609.09841v2 (CC BY 4.0)}
\begin{document}
\maketitle

\noindent\emph{Source and licence.} \bf A Sharp Sobolev Inequality on Compact CR Manifolds. Version arXiv:2609.09841v2 (CC BY 4.0), distributed by arXiv under the Creative Commons Attribution 4.0 International licence, http://creativecommons.org/licenses/by/4.0/, as recorded in the arXiv licence metadata for this exact version and shown on the abstract page https://arxiv.org/abs/2609.09841v2. Full licence notice: \texttt{licenses/arxiv-2609.09841-CC-BY-4.0.txt}.\par\smallskip

\noindent\textit{Editorial scope.} Counted unit: the article's lemma lemmapohozaev (source0.tex lines 988--1000). The article's complete original proof of it is delivered with it (source0.tex lines 1002--1126, one \texttt{proof} environment of 125 lines). No mathematics is written, repaired or re-derived: the statement and the proof are the article's own lines, in the article's own notation, and no retained line is substituted or re-flowed.  Retained uncounted as the source-local context the proof needs: lines 430--442; lines 446--450; lines 533--535; lines 537--539; lines 899--904; lines 980--982.  Nothing was omitted from any retained span, so the package carries no omission record.  No retained line contains a citation, so the wrapper carries no bibliography and no bibliography entry.  No retained line contains a figure environment, an image-inclusion command or drawing code, so no image is delivered and the package declares no asset dependency.  Editorial formatting only, in addition: an AMS \texttt{amsart} preamble that restates the article's own declarations and macros used by the retained lines, namely the environments \texttt{corollary}, \texttt{lemma} on separate counters, together with the standard packages those lines need.  The retained environments print in this document as Corollary 1, Lemma 1 in order of appearance, whereas the article prints the same items with the numbering of its own full document.  The complete native source of the article as distributed in the arXiv e-print for this exact version is bundled unchanged as \texttt{sources/209855/source0.tex}.
\begin{equation}\label{eq:JL-frame-comparison}
\begin{aligned}
W_\alpha
&=
Z_\alpha
+
\sum_{\beta=1}^n O(\rho_0^2)Z_\beta
+
\sum_{\beta=1}^n O(\rho_0^2)Z_{\bar\beta}
+
O(\rho_0^3)T .
\end{aligned}
\end{equation}

\begin{equation}\label{eq:JL-volume-comparison}
\Psi_q^*dv_\theta
=
(1+O(\rho_0^2))\,dv_0 .
\end{equation}

		\begin{equation}\label{eqnmin}
			I_\ell(u_\ell)=\lambda_\ell:=\inf_{u\in S^{1,2}(M)\setminus\{0\}} I_\ell(u)<K(n)^{-2},\quad\|u_\ell\|_{L^{Q^*}(M)}=1.
		\end{equation}

		\begin{equation}\label{eqnEL}
			-\Delta_b u_\ell+\ell u_\ell=\lambda_\ell u_\ell^{Q^*-1}, \quad u_\ell > 0, \quad \text { on } M .
		\end{equation}

	\begin{corollary}\label{estimatering}
		For every \(0 < r \le r_G\), there exists  \(\ell_r>0\) and a constant \(C>0\), independent of \(\ell\), such that for all \(\ell>\ell_r\),
		\[
		\int_{B_{r}(q_\ell)} u_\ell^2\,dv_\theta \le C\int_{B_{r/2}(q_\ell)} u_\ell^2\,dv_\theta,
		\]
	\end{corollary}

\begin{equation}\label{eq:cutoff-mass-lower}
\int w_\ell^{Q^*}\,dv_\theta\ge c_0
\end{equation}

\begin{lemma}\label{lemmapohozaev}
For all sufficiently large \(\ell\),
\[
\ell\int w_\ell^2\,dv_\theta
\le
C\int \rho^2w_\ell^{Q^*}\,dv_\theta
+
C\int \rho^2|\nabla_bw_\ell|^2\,dv_\theta 
+
C\int \rho^4|T_\theta w_\ell|^2\,dv_\theta,
\]
where \(C>0\) is independent of \(\ell\).
\end{lemma}

\begin{proof}
Multiplying equation \eqref{eqnEL}
by \(\eta_\ell^2u_\ell\) and integrating by parts, we obtain
\begin{equation}\label{eq:cutoff-energy-identity}
\int|\nabla_bw_\ell|^2\,dv_\theta
+
\ell\int w_\ell^2\,dv_\theta=
\lambda_\ell
\int\eta_\ell^2u_\ell^{Q^*}\,dv_\theta \\
+
\int|\nabla_b\eta_\ell|^2u_\ell^2\,dv_\theta.
\end{equation}

By H\"older's inequality and the normalization of \(u_\ell\) in \eqref{eqnmin},
\[
\int\eta_\ell^2u_\ell^{Q^*}\,dv_\theta
\le
\Big(\int w_\ell^{Q^*}\,dv_\theta\Big)^{2/Q^*}
\Big(\int_Mu_\ell^{Q^*}\,dv_\theta\Big)^{(Q^*-2)/Q^*}
=
\Big(\int w_\ell^{Q^*}\,dv_\theta\Big)^{2/Q^*}.
\]
Moreover, using the uniform bound for \(\nabla_b\eta_\ell\) and applying Corollary~\ref{estimatering} with the fixed radius \(r_P\), we have
\[
\int|\nabla_b\eta_\ell|^2u_\ell^2\,dv_\theta
\le
C\int_{B_{r_P}(q_\ell)}u_\ell^2\,dv_\theta 
\le
C\int_{B_{r_P/2}(q_\ell)}u_\ell^2\,dv_\theta 
\le
C\int w_\ell^2\,dv_\theta \le \frac{\ell}{2}\int w_\ell^2\,dv_\theta
\]
when \(\ell\) is sufficiently large.
It follows from \eqref{eq:cutoff-energy-identity} that
\begin{equation}\label{eq:cutoff-below-sharp}
\int|\nabla_bw_\ell|^2\,dv_\theta
+
\frac{\ell}{2}\int w_\ell^2\,dv_\theta
\le
\lambda_\ell
\Big(\int w_\ell^{Q^*}\,dv_\theta\Big)^{2/Q^*}\le K(n)^{-2}
\Big(\int w_\ell^{Q^*}\,dv_\theta\Big)^{2/Q^*}.
\end{equation}

We now compare the right-hand side of
\eqref{eq:cutoff-below-sharp} with the sharp inequality on the Heisenberg
group.  Define
\[
\widetilde w_\ell(\xi)
:=
w_\ell(\Psi_\ell(\xi)),
\qquad
\xi\in B_{r_P}(0),
\]
and extend \(\widetilde w_\ell\) by zero outside \(B_{r_P}(0)\).
Since \(w_\ell\) vanishes near \(\partial B_{r_P}(q_\ell)\), this extension
belongs to \(S^{1,2}(\mathbb H^n)\).


By the inverse form of the Jerison--Lee frame expansion \eqref{eq:JL-frame-comparison}, together with
Young's inequality, we have
\[
|\nabla_H\widetilde w_\ell(\xi)|^2
\le
(1+C\rho(\Psi_\ell(\xi))^2)
|\nabla_bw_\ell|^2(\Psi_\ell(\xi))
+
C\rho(\Psi_\ell(\xi))^4
|T_\theta w_\ell|^2(\Psi_\ell(\xi)).
\]
Moreover, the inverse form of \eqref{eq:JL-volume-comparison} gives
\(
dv_0=(1+O(\rho^2))\Psi_\ell^*dv_\theta
\)
and therefore yields
\begin{equation}\label{eqexpnablab}
\int_{\mathbb H^n}|\nabla_H\widetilde w_\ell|^2\,dv_0
\le
\int |\nabla_bw_\ell|^2\,dv_\theta
+
C\int \rho^2|\nabla_bw_\ell|^2\,dv_\theta  
+
C\int \rho^4|T_\theta w_\ell|^2\,dv_\theta .
\end{equation}
By the inverse volume comparison and
\eqref{eq:cutoff-mass-lower}, the mean value theorem gives
\begin{equation}\label{eqexpqstar}
\Big(\int w_\ell^{Q^*}\,dv_\theta\Big)^{2/Q^*}
\le
\Big(\int_{\mathbb H^n}
\widetilde w_\ell^{Q^*}\,dv_0\Big)^{2/Q^*}+
C\int\rho^2w_\ell^{Q^*}\,dv_\theta.
\end{equation}
Combining \eqref{eq:cutoff-below-sharp},
\eqref{eqexpnablab}, and \eqref{eqexpqstar} with the sharp
Folland--Stein inequality on \(\mathbb H^n\), we obtain
\[
\begin{aligned}
\int |\nabla_bw_\ell|^2\,dv_\theta
+\frac{\ell}{2}\int w_\ell^2\,dv_\theta
&\le
K(n)^{-2}
\Big(\int w_\ell^{Q^*}\,dv_\theta\Big)^{2/Q^*} \\
&\le
K(n)^{-2}
\Big(\int_{\mathbb H^n}\widetilde w_\ell^{Q^*}\,dv_0\Big)^{2/Q^*}
+
C\int\rho^2w_\ell^{Q^*}\,dv_\theta \\
&\le
\int_{\mathbb H^n}|\nabla_H\widetilde w_\ell|^2\,dv_0
+
C\int\rho^2w_\ell^{Q^*}\,dv_\theta \\
&\le
\int |\nabla_bw_\ell|^2\,dv_\theta
+
C\int\rho^2|\nabla_bw_\ell|^2\,dv_\theta \\
&\quad
+
C\int\rho^4|T_\theta w_\ell|^2\,dv_\theta
+
C\int\rho^2w_\ell^{Q^*}\,dv_\theta .
\end{aligned}
\]
This proves the lemma.
\end{proof}

\end{document}
